% Mathématiques 4e — cours V3
% Algorithmique : variables, conditions et boucles

\section{Objectifs}

\begin{itemize}
\item Comprendre le rôle d'une variable informatique.
\item Lire et modifier la valeur d'une variable.
\item Distinguer entrée, traitement et sortie.
\item Traduire une formule mathématique en instructions.
\item Utiliser une instruction conditionnelle.
\item Comprendre la structure « si ... alors ... sinon ».
\item Utiliser une boucle à nombre de répétitions fixé.
\item Simuler un calcul ou une expérience simple.
\item Analyser et corriger un programme.
\item Prévoir le résultat d'un programme avant de l'exécuter.
\end{itemize}

\section{1. Programme et algorithme}

Un algorithme décrit une suite d'étapes permettant de réaliser une tâche.

Un programme traduit cet algorithme dans un langage compris par une machine.

En mathématiques, on cherche surtout à comprendre :
\begin{itemize}
\item les données d'entrée ;
\item les calculs ;
\item les décisions ;
\item les répétitions ;
\item le résultat produit.
\end{itemize}

\section{2. Variable informatique}

\begin{definition}
Une variable informatique est une zone de mémoire nommée qui contient une valeur pouvant être lue ou modifiée pendant l'exécution.
\end{definition}

\coursefigure{/images/mathematiques/4eme/algorithmique-variable.svg}{Une variable nommée score contient une valeur puis est modifiée par une instruction.}{La variable peut changer de valeur pendant l'exécution : elle mémorise l'état courant du calcul.}

Si :
\[
\text{score}=7,
\]
puis que l'on exécute :
\[
\text{score}\leftarrow\text{score}+3,
\]
la nouvelle valeur est :
\[
10.
\]

\section{3. Affectation}

L'instruction :
\[
x\leftarrow5
\]
signifie :
\[
\text{« mettre la valeur }5\text{ dans la variable }x\text{ »}.
\]

Elle ne doit pas être confondue avec l'égalité mathématique :
\[
x=5.
\]

Dans un programme, la valeur de $x$ pourra ensuite changer.

\section{4. Lire puis modifier}

Considérons :
\begin{verbatim}
x prend la valeur 4
x prend la valeur x + 3
x prend la valeur 2 * x
\end{verbatim}

Après la première instruction :
\[
x=4.
\]

Puis :
\[
x=7.
\]

Enfin :
\[
x=14.
\]

La sortie finale vaut :
\[
\boxed{14}.
\]

\section{5. Entrée et sortie}

Une entrée est une donnée fournie au programme.

Une sortie est une information produite.

Exemple :
\begin{verbatim}
demander un nombre
mettre la réponse dans x
calculer 3*x+2
afficher le résultat
\end{verbatim}

Si l'entrée vaut :
\[
5,
\]
la sortie vaut :
\[
3\times5+2=17.
\]

\section{6. Traduire une formule}

Pour :
\[
A=\pi r^2,
\]
un programme peut exécuter :
\begin{verbatim}
lire r
A = pi * r * r
afficher A
\end{verbatim}

Les priorités opératoires doivent être respectées.

\section{7. Condition}

Une condition est une affirmation qui peut être vraie ou fausse.

Exemples :
\[
x>10,
\]
\[
n=0,
\]
\[
\text{score}\ge20.
\]

Le programme peut choisir une action en fonction du résultat du test.

\coursefigure{/images/mathematiques/4eme/algorithmique-condition.svg}{Organigramme simple montrant une condition avec une branche oui et une branche non.}{Une instruction conditionnelle permet au programme de prendre des décisions différentes selon que le test est vrai ou faux.}

\section{8. Structure « si ... alors »}

\begin{verbatim}
si x > 10 alors
    afficher "grand"
fin si
\end{verbatim}

Si :
\[
x=14,
\]
le message est affiché.

Si :
\[
x=8,
\]
aucune action n'est exécutée dans le bloc.

\section{9. Structure « si ... alors ... sinon »}

\begin{verbatim}
si x >= 0 alors
    afficher "positif ou nul"
sinon
    afficher "négatif"
fin si
\end{verbatim}

Une seule des deux branches est exécutée.

\section{10. Exemple : parité}

Un entier $n$ est pair si son reste dans la division par $2$ est nul.

On peut écrire :
\begin{verbatim}
si n modulo 2 = 0 alors
    afficher "pair"
sinon
    afficher "impair"
fin si
\end{verbatim}

Pour :
\[
n=18,
\]
la condition est vraie.

Pour :
\[
n=23,
\]
elle est fausse.

\section{11. Comparer deux nombres}

Programme :
\begin{verbatim}
lire a
lire b
si a > b alors
    afficher a
sinon
    afficher b
fin si
\end{verbatim}

Ce programme affiche le plus grand des deux nombres lorsque :
\[
a\neq b.
\]

Si :
\[
a=b,
\]
il affiche $b$, qui est égal à $a$.

\section{12. Condition composée}

On peut parfois combiner deux tests.

Par exemple, vérifier que :
\[
0\le x\le10.
\]

Dans un langage par blocs, on peut tester :
\[
x\ge0
\]
et :
\[
x\le10.
\]

Les deux conditions doivent être vraies.

\section{13. Boucle inconditionnelle}

Une boucle répète une suite d'instructions.

Exemple :
\begin{verbatim}
répéter 5 fois
    avancer de 40
    tourner de 72 degrés
\end{verbatim}

Cette séquence peut construire un pentagone régulier.

Le nombre de répétitions est fixé avant l'exécution.

\section{14. Variable accumulatrice}

Une variable peut accumuler une somme.

\begin{verbatim}
somme = 0
répéter 4 fois
    lire un nombre
    somme = somme + nombre
afficher somme
\end{verbatim}

Après chaque passage, la variable :
\[
\text{somme}
\]
mémorise le total provisoire.

\section{15. Compteur}

Un compteur augmente d'une unité à chaque événement.

\begin{verbatim}
compteur = 0
répéter 100 fois
    simuler un lancer
    si succès alors
        compteur = compteur + 1
\end{verbatim}

À la fin :
\[
\text{compteur}
\]
est le nombre de succès.

\section{16. Simulation d'une fréquence}

Après :
\[
n
\]
essais, si le compteur vaut :
\[
s,
\]
la fréquence expérimentale est :
\[
f=\dfrac sn.
\]

Un programme permet de répéter rapidement une expérience aléatoire simple.

\section{17. Table de suivi}

Pour comprendre un programme, on peut suivre les variables étape par étape.

Programme :
\begin{verbatim}
x = 3
y = 2*x + 1
x = y - 4
\end{verbatim}

Table de suivi :

\begin{tabular}{c|c|c}
Étape & $x$ & $y$ \\
initialisation & $3$ & -- \\
après calcul de $y$ & $3$ & $7$ \\
après modification de $x$ & $3$ & $7$ \\
\end{tabular}

La dernière instruction donne en réalité :
\[
x=7-4=3.
\]

La valeur finale de $x$ reste donc :
\[
3.
\]

\section{18. Prévoir avant d'exécuter}

Considérons :
\begin{verbatim}
x = 6
si x > 4 alors
    x = 2*x
sinon
    x = x+10
afficher x
\end{verbatim}

La condition :
\[
6>4
\]
est vraie.

On exécute donc :
\[
x=2\times6=12.
\]

La sortie vaut :
\[
\boxed{12}.
\]

\section{19. Déboguer}

Déboguer signifie rechercher et corriger une erreur.

\begin{methode}
\begin{enumerate}
\item prévoir le résultat attendu ;
\item exécuter ou simuler pas à pas ;
\item suivre la valeur des variables ;
\item repérer la première étape qui diverge ;
\item corriger une seule erreur à la fois ;
\item tester à nouveau.
\end{enumerate}
\end{methode}

\section{20. Erreur de condition}

On souhaite afficher « admis » si :
\[
\text{note}\ge10.
\]

Le programme contient :
\begin{verbatim}
si note > 10 alors
    afficher "admis"
\end{verbatim}

La note :
\[
10
\]
n'est alors pas reconnue comme admise.

Il faut écrire :
\[
\text{note}\ge10.
\]

\section{21. Erreur d'affectation}

Supposons :
\begin{verbatim}
x = 5
x = x + 1
x = x + 1
\end{verbatim}

La valeur finale est :
\[
7.
\]

Chaque affectation remplace la valeur précédente.

On ne peut pas lire :
\[
x=x+1
\]
comme une égalité mathématique ordinaire.

\section{22. Programme de calcul}

On choisit un nombre $x$.

Le programme :
\begin{verbatim}
y = x*x
y = y - 4
\end{verbatim}

traduit :
\[
y=x^2-4.
\]

Pour :
\[
x=-3,
\]
on obtient :
\[
y=9-4=5.
\]

\section{23. Programme et équation}

Si l'on cherche une entrée produisant :
\[
y=21
\]
dans le programme :
\[
y=4x+1,
\]
on résout :
\[
4x+1=21.
\]

Donc :
\[
x=5.
\]

L'algorithmique et le calcul littéral décrivent ici le même modèle.

\section{24. Géométrie dynamique}

Un programme peut déplacer un point ou un objet selon des instructions.

Exemple :
\begin{verbatim}
répéter 4 fois
    avancer de 80
    tourner de 90 degrés
\end{verbatim}

Le tracé forme un carré.

Modifier :
\[
80
\]
change la taille.

Modifier :
\[
90^\circ
\]
change la forme du tracé.

\section{25. Méthode d'analyse}

\begin{methode}
Pour analyser un programme :
\begin{enumerate}
\item relever les variables ;
\item repérer leurs valeurs initiales ;
\item suivre les affectations dans l'ordre ;
\item identifier les conditions ;
\item identifier les boucles ;
\item prévoir la sortie ;
\item vérifier par une table de suivi si nécessaire.
\end{enumerate}
\end{methode}

\section{26. Erreurs fréquentes}

\begin{itemize}
\item Confondre affectation et égalité mathématique.
\item Oublier que la valeur d'une variable peut changer.
\item Exécuter les deux branches d'une condition.
\item Inverser les comparateurs $>$ et $<$.
\item Oublier le cas d'égalité dans une condition.
\item Confondre compteur et valeur simulée.
\item Oublier d'initialiser une variable accumulatrice.
\item Lire une boucle comme si son contenu n'était exécuté qu'une seule fois.
\end{itemize}

\section{À retenir}

\begin{aretenir}
Une variable mémorise une valeur susceptible d'être modifiée :
\[
\boxed{x\leftarrow\text{nouvelle valeur}}.
\]

Une condition permet de choisir une action :
\[
\boxed{\text{si condition alors ... sinon ...}}.
\]

Une boucle répète une séquence d'instructions.

Pour comprendre ou corriger un programme, il faut suivre l'évolution des variables étape par étape et prévoir le résultat avant l'exécution.
\end{aretenir}
